% unidades/20-cuerpo-salud.tex
\chapter{🏥 体と健康 — Cuerpo y Salud}
\unidadheader{20}{体と健康}{Cuerpo y Salud}{Hablar de síntomas y en el médico}

\begin{tcolorbox}[vocabbox, title={📌 Vocabulario Nuevo}]
\begin{tabularx}{\linewidth}{llX}
\toprule
\textbf{Japonés} & \textbf{Español} & \textbf{Tipo} \\
\midrule
\vocabitem{頭}{あたま}{cabeza}{sust.}
\vocabitem{目}{め}{ojo}{sust.}
\vocabitem{耳}{みみ}{oreja / oído}{sust.}
\vocabitem{歯}{は}{diente}{sust.}
\vocabitem{足}{あし}{pie / pierna}{sust.}
\vocabitem{お腹}{おなか}{estómago / barriga}{sust.}
\vocabitem{熱}{ねつ}{fiebre}{sust.}
\vocabitem{風邪}{かぜ}{resfriado}{sust.}
\vocabitem{薬}{くすり}{medicina}{sust.}
\vocabitem{医者}{いしゃ}{médico/a}{sust.}
\bottomrule
\end{tabularx}
\end{tcolorbox}

\subsection*{🗣️ Frases Clave}

\frase{\ruby{頭}{あたま}が\ruby{痛}{いた}いです。}{Me duele la cabeza.}
\frase{\ruby{熱}{ねつ}があります。}{Tengo fiebre.}
\frase{\ruby{風邪}{かぜ}を\ruby{引}{ひ}きました。}{Me resfrié.}
\frase{\ruby{薬}{くすり}を\ruby{飲}{の}んでください。}{Tome la medicina.}
\frase{\ruby{今日}{きょう}は\ruby{無理}{むり}しないでください。}
  {No se esfuerce de más hoy.}

\subsection*{⚙️ Gramática}

\patron{Expresar dolor}
  {部位 + が + 痛いです}
  {\ej{\ruby{歯}{は}が\ruby{痛}{いた}いです。}{Me duele un diente.}
  \ej{のどが\ruby{痛}{いた}いです。}{Me duele la garganta.}}

\patron{Prohibición suave: ないでください}
  {V-ない + でください (por favor, no haga ~)}
  {\ej{\ruby{無理}{むり}しないでください。}
    {Por favor, no se esfuerce de más.}
  \ej{\ruby{辛}{から}いものを\ruby{食}{た}べないでください。}
    {Por favor, no coma cosas picantes.}}

\patron{Consejo: ないほうがいいです}
  {V-ない + ほうがいいです (es mejor no hacer ~)}
  {\ej{お\ruby{酒}{さけ}を\ruby{飲}{の}まないほうがいいです。}
    {Es mejor que no deba alcohol.}
  \ej{しばらく\ruby{運動}{うんどう}しないほうがいいです。}
    {Es mejor que no haga ejercicio por un tiempo.}}

\begin{tcolorbox}[ejerciciobox, title={📝 Ejercicios Rápidos}]
\begin{enumerate}
  \item Di tres síntomas en japonés.
  \item Usa ないほうがいいです para dar un consejo a alguien con fiebre.
  \item Traduce: \ruby{病院}{びょういん}に\ruby{行}{い}ったほうがいいですよ。
\end{enumerate}
\vspace{0.4em}
\textbf{Respuestas:}
1) (ej: \ruby{熱}{ねつ}があります、せきが\ruby{出}{で}ます、のどが\ruby{痛}{いた}いです)\quad
2) お\ruby{風呂}{ふろ}に\ruby{入}{はい}らないほうがいいです。\quad
3) Es mejor que vayas al hospital.
\end{tcolorbox}

\begin{tcolorbox}[kanjibox, title={🀄 Kanjis de la Unidad}]
\begin{tabularx}{\linewidth}{cllXX}
\toprule
\textbf{Kanji} & \textbf{On} & \textbf{Kun} & \textbf{Significado} & \textbf{Vocab} \\
\midrule
\kanjirow{体}{たい}{からだ}{cuerpo}{\ruby{体}{からだ} / \ruby{体育}{たいいく}}
\kanjirow{健}{けん}{—}{salud / sano}{\ruby{健康}{けんこう} / \ruby{健全}{けんぜん}}
\kanjirow{康}{こう}{—}{salud / paz}{\ruby{健康}{けんこう}}
\kanjirow{痛}{つう}{いた}{dolor}{\ruby{痛}{いた}い / \ruby{頭痛}{ずつう}}
\kanjirow{熱}{ねつ}{あつ}{fiebre / calor}{\ruby{熱}{ねつ} / \ruby{熱}{あつ}い}
\bottomrule
\end{tabularx}
\end{tcolorbox}
